\section{Method}
\label{sec:method}

\subsection{Variational Free Energy}
\label{subsec:first-principles}

The robot cannot directly observe its true field pose. It must infer a hidden state from sensory observations under a generative model. In Bayesian terms, the goal is to estimate a posterior belief over hidden state $x$ given observations $y$:
\begin{equation}
    p(x \mid y) =
    \frac{p(y \mid x)p(x)}{p(y)}.
    \label{eq:bayes}
\end{equation}
This posterior is generally intractable. Instead, a tractable belief $q(x;\zeta)$ is optimized so that it approximates the posterior over hidden state $x$. In active inference, this optimization is expressed as minimization of variational free energy, see Appendix~\ref{app:vfe-derivation}:
\begin{equation}
    \underset{\zeta}{\arg\min}\, \mathcal{F}(y,\zeta)
    =
    \underset{\zeta}{\arg\min}
    \int q(x;\zeta)
    \ln\left(
        \frac{q(x;\zeta)}{p(x,y)}
    \right) dx .
    \label{eq:vfe-general}
\end{equation}
The free-energy objective can be evaluated as a function of sensory data and belief-state parameters. It can therefore be minimized as an automatable optimization problem. Minimizing $\mathcal{F}(y,\zeta)$ with respect to $\zeta$ optimizes $q(x;\zeta)$ so that it approximates the posterior over hidden state $x$ given sensory observations $y$. For a Gaussian belief, as used in the Laplace approximation below, $\zeta$ corresponds to the mean and variance.
\subsection{Laplace-Encoded Free Energy}
\label{subsec:gaussian-errors}

Under the Laplace and mean-field approximations, the variational free energy in Eq.~\ref{eq:vfe-general} reduces to a sum of squared precision-weighted prediction errors; see Appendix~\ref{app:laplace-free-energy}. With belief parameters $\mu$, prediction error vector $\epsilon$, and precision matrix $\Pi$ (inverse variance), this gives
\begin{equation}
    \underset{\tilde{\mu}}{\arg\min}\, \mathcal{F}(y,\tilde{\mu})
    =
    \underset{\tilde{\mu}}{\arg\min}
    \left(
        \frac{1}{2}\tilde{\epsilon}^\top \tilde{\Pi} \tilde{\epsilon}
        -
        \frac{1}{2}\ln|\tilde{\Pi}|
    \right).
    \label{eq:vfe-gaussian}
\end{equation}
The tilde denotes generalized-coordinate quantities; in this paper, each pose variable is represented as a vector such as $\tilde{x}=[x,\dot{x},\ddot{x}]^\top$ for position, velocity, and acceleration. The compact formulation in Eq.~\ref{eq:vfe-gaussian} can be unpacked for our non-hierarchical robot-soccer model as:

\begin{align}
    \underset{\mu}{\arg\min}\, \mathcal{F}(y,\mu)
    &=
    \underset{\mu}{\arg\min}
    \Bigg[
    \frac{1}{2}
    \sum_{j=1}^{P+1}
    \sum_{k=1}^{n}
    \left(
        \Pi_{g(i,j,k)}
        \epsilon^2_{g(i,j,k)}
        -
        \ln\left(\Pi_{g(i,j,k)}\right)
    \right)
    \nonumber \\
    &\qquad\qquad+
    \frac{1}{2}
    \sum_{j=1}^{P+1}
    \sum_{k=1}^{q}
    \left(
        \Pi_{p(i,j,k)}
        \epsilon^2_{p(i,j,k)}
        -
        \ln\left(\Pi_{p(i,j,k)}\right)
    \right)
    \Bigg].
    \label{eq:vfe-motion-place-expanded}
\end{align}

The first term contains the motion prediction errors $g$. The hidden state is one robot field-pose belief with three pose variables, $x$, $y$, and $\theta$, so $n=3$. Each pose variable is represented up to second order in generalized coordinates, namely position, velocity, and acceleration, so $P=2$. This gives 9 motion prediction errors to be minimized.

The second term contains the sensory landmark prediction errors $p$. In this paper the landmarks are static (only position, no speed or acceleration), and up to 9 landmarks are used, giving 27 precision-weighted landmark prediction errors to be minimized.


\medskip
Free energy is therefore low when one field-pose belief is jointly consistent with the robot's motion evidence and its observed field landmarks. Minimizing the free energy means finding the three generalized-coordinate belief vectors $\tilde{\mu}_{g,x}$, $\tilde{\mu}_{g,y}$, and $\tilde{\mu}_{g,\theta}$, or nine scalar belief parameters in total, that minimize these 36 prediction errors. This minimization is done by gradient descent.

\subsection{Gradient Descent on the Free Energy}
\label{subsec:gradient-descent-free-energy}

 Taking the free-energy gradient with respect to $\tilde{\mu}_{g}$ gives the generalized-coordinate belief update, see Appendix~\ref{app:slide15-update-derivation}:
\begin{equation}
    \dot{\tilde{\mu}}_{g}
    =
    D\tilde{\mu}_{g}
    +
    \underbrace{
        \left[
            -
            D^\top \tilde{\Pi}_{g}\tilde{\epsilon}_{g}
            +
            \left(
                \partial_{\tilde\mu_g}\tilde{f}(\tilde{\mu}_{g})
            \right)^\top
            \tilde{\Pi}_{g}\tilde{\epsilon}_{g}
        \right]
    }_{\text{motion prediction errors}}
    +
    \underbrace{
        \sum_{\mathrm{landmarks}}
        \left(
            \partial_{\tilde\mu_g}\tilde{g}(\tilde{\mu}_{g})
        \right)^\top
        \tilde{\Pi}_{p}\tilde{\epsilon}_{p}
    }_{\text{landmark prediction errors}}.
    \label{eq:slide15-generic-update}
\end{equation}
with prediction errors
\begin{equation}
    \tilde{\epsilon}_g
    =
    D\tilde{\mu}_g
    -
    \tilde{f}(\tilde{\mu}_g),
    \qquad
    \tilde{\epsilon}_p
    =
    \tilde{y}
    -
    \tilde{g}(\tilde{\mu}_g).
    \label{eq:prediction-errors}
\end{equation}

\subsubsection{Motion Part}
\label{subsec:gradient-motion}

The motion part of Eq.~\ref{eq:slide15-generic-update} is the motion prediction-error contribution.
To keep the notation compact, the following equations show the $x$ component, $\tilde{\mu}_{g,x}=[\mu_{g,x},\mu'_{g,x},\mu''_{g,x}]^\top$, which is the robot's generalized-coordinate belief for its field $x$ position in the Field Coordinate System (FCS). Likewise for $y$ and $\theta$.
\begin{align}
    \tilde{\epsilon}_{g,x}
    &=
    D\tilde{\mu}_{g,x}
    -
    \tilde{f}(\tilde{\mu}_g)_x
    =
    \begin{bmatrix}
        \mu'_{g,x} \\
        \mu''_{g,x} \\
        0
    \end{bmatrix}
    -
    \begin{bmatrix}
        f_{\mathrm{RCS}\rightarrow\mathrm{FCS}}(v^{\mathrm{odom}})_x \\
        f_{\mathrm{RCS}\rightarrow\mathrm{FCS}}(a^{\mathrm{IMU}})_x \\
        0
    \end{bmatrix}
    \nonumber \\
    &=
    \text{motion belief in FCS}
    -
    \text{predicted motion in FCS}.
    \label{eq:motion-error-x}
\end{align}
Here, $D$ is the generalized-coordinate shift operator: it shifts the vector entries up by one and adds a zero at the bottom.
The term $\tilde{f}$ is the motion generative model. In this implementation, we choose to use the available sensory odometry and IMU measurements directly, rather than motor commands or a separate movement function. To make use of these motion sensory observations, they are translated from RCS into FCS. Since the soccer-field geometry is known, this is a straightforward RCS-to-FCS translation using the belief heading, see Appendix~\ref{app:coordinate-conventions}, and the same transform is used for both velocity and acceleration observations.
The corresponding motion precision matrix for the $x$ component is
\begin{equation}
    \tilde{\Pi}_{g,x}
    =
    \begin{bmatrix}
        \Pi_{\mathrm{odo}} & 0 & 0 \\
        0 & \Pi_{\mathrm{IMU}} & 0 \\
        0 & 0 & 0
    \end{bmatrix},
    \label{eq:motion-precision-x}
\end{equation}
which contains the precisions of the odometry and IMU prediction errors. In this implementation these precisions are fixed parameters and set to 1, and they can be tuned if required. Motion in $x$, $y$, and $\theta$ is treated as factorized, matching the robot-soccer setting where omnidirectional wheels or humanoid stepping can produce independent motion components. As a design choice, the term $\left(\partial_{\tilde\mu_g}\tilde{f}(\tilde{\mu}_g)\right)^\top$ is zero because the motion is not dependent on the current field position or angle.

\subsubsection{Landmark Part}
\label{subsec:gradient-landmarks}

The landmark part of Eq.~\ref{eq:slide15-generic-update} is the sensory prediction-error contribution.
To keep the notation compact, the following equations show the x coordinate of one unambiguous observed landmark example.
\begin{align}
    \tilde{\epsilon}_{p,x}
    &=
    \begin{bmatrix}
        y_x \\
        0 \\
        0
    \end{bmatrix}
    -
    \begin{bmatrix}
        g_{\mathrm{FCS}\rightarrow\mathrm{RCS}}(\mu_g)_x \\
        0 \\
        0
    \end{bmatrix}
    \nonumber \\
    &=
    \text{observed landmark in RCS}
    -
    \text{predicted landmark in RCS}.
    \label{eq:landmark-error-x}
\end{align}
Here, $y_x$ is the observed landmark $x$ position in RCS and $g_{\mathrm{FCS}\rightarrow\mathrm{RCS}}(\mu_g)_x$ is the predicted landmark $x$ position in RCS. The same construction is used for the observed landmark $y$ position and $\theta$ angle.
Landmark observations provide relative $x$ and $y$ positions, not relative velocity or acceleration, so there are no generalized-coordinate rows for landmark velocity or acceleration in the sensory prediction error.
The soccer field geometry is predefined, so the landmark observation model and its Jacobian can be calculated straightforwardly from the known FCS/RCS coordinate conventions, see Appendix~\ref{app:coordinate-conventions}.
The term $\partial_{\tilde\mu_g}\tilde{g}(\tilde{\mu}_g)$ in Eq.~\ref{eq:slide15-generic-update} is the landmark Jacobian $\boldsymbol{J}_{\ell}$. It describes how small changes in the robot pose, including $x$, $y$, and $\theta$, affect the predicted landmark observation.

For the observed landmark $x$ coordinate, the corresponding landmark precision matrix is
\begin{equation}
    \tilde{\Pi}_{p,x}
    =
    \begin{bmatrix}
        \Pi_{\mathrm{cam}} & 0 & 0 \\
        0 & 0 & 0 \\
        0 & 0 & 0
    \end{bmatrix}.
    \label{eq:landmark-precision-x}
\end{equation}

\subsection{Distance-Weighted Precision Modifier}
\label{subsec:distance-weighted-precision}

A useful optional refinement is to make landmark precision depend on observation distance. Camera-based position estimates generally become less precise at longer distances because the pixel-per-meter ratio decreases with range. Distance weighting provides a simple way to represent this range-dependent confidence in the landmark update. It is not required for convergence: suitable learning-rate settings can also yield stable behavior. As a practical design choice, we use a precision distance modifier based on the landmark's predicted position in RCS:

\begin{equation}
d_\ell =
\sqrt{
    \left(g(\mu_g)_{x,\ell}\right)^2
    +
    \left(g(\mu_g)_{y,\ell}\right)^2
},
    \qquad
    \rho_\ell
    =
    \max
    \left(
        \rho_{\min},
        \frac{1}{1 + \left(d_\ell / d_{\mathrm{ref}}\right)^2}
    \right).
    \label{eq:distance-precision-modifier}
\end{equation}
In the proof-of-concept experiments, $d_{\mathrm{ref}}=600$ cm and $\rho_{\min}=0.15$.

\subsection{Resulting Update Rules}
\label{subsec:update-rules}

For the $x$ field-pose component, this gives the following straightforward
matrix calculation. Here, $\boldsymbol{J}_{g_x,\ell}$ is the column of the landmark Jacobian associated with belief component $g_x$, and the scalar factors $\eta^{c}_{xy}$, $\eta^{v}_{xy}$, and $\eta^{a}_{xy}$ denote learning rates for the camera, odometry-velocity, and IMU-acceleration update terms.

{
\small
\begin{align}
\dot{\tilde{\boldsymbol{\mu}}}_{g_x}
&\!=\!
\begin{bmatrix}
\mu'_{g_x}\\
\mu''_{g_x}\\
0
\end{bmatrix}
\!-\!
\begin{bmatrix}
0\\
\eta^{v}_{xy}\Pi_{\mathrm{odo}}
\left(\mu'_{g_x}-f(v^{\mathrm{odom}})_x\right)\\
\eta^{a}_{xy}\Pi_{\mathrm{IMU}}
\left(\mu''_{g_x}-f(a^{\mathrm{IMU}})_x\right)
\end{bmatrix}
\!+\!
\sum_{\ell}
\begin{bmatrix}
\eta^{c}_{xy}
\rho_\ell \boldsymbol{J}_{g_x,\ell}^{\top}
\boldsymbol{\Pi}_{\mathrm{cam}}
\left(y_{x,\ell}-g(\mu_g)_{x,\ell}\right)\\
0\\
0
\end{bmatrix}.
\end{align}
}
The same calculation is applied to the $y$ and $\theta$ belief components using their corresponding landmark-Jacobian columns and learning rates. The resulting generalized-coordinate update is integrated with Euler's forward method using a configurable step size corresponding to the frequency of the robot localization loop. The learning rates are manually tuned. The update has a simple intuition. First, the motions of the current generalized coordinates are shifted forward in time. Then the motion terms correct the belief: if the predicted movement is too fast, the acceleration is reduced. Next, the landmark terms correct the belief: If the position is too far, the velocity is reduced.

 The resulting update rules are efficient matrix calculations over $x$, $y$, and $\theta$, suitable for a high-frequency robot-soccer localization loop.
